\documentclass[12pt,a4paper]{article}

\usepackage[utf8]{inputenc}
\usepackage[T1]{fontenc}
\usepackage{amsmath}
\usepackage{geometry}
\usepackage{booktabs}
\usepackage{float}

\geometry{margin=1in}

\title{Report 6: Convolutional Neural Network}
\author{Dao Quy Tung}
\date{Deep Learning 2026}

\begin{document}

\maketitle

\section{Introduction}

The objective of this labwork is to implement a convolutional neural network for image classification.
The selected architecture follows the VGG19 block structure and is implemented using PyTorch.
The model is trained from scratch, so no pretrained weights are loaded.

The experiment uses a generated image dataset with two classes:

\begin{table}[H]
\centering
\begin{tabular}{ll}
\toprule
Class & Description \\
\midrule
\texttt{horizontal} & Image containing a horizontal bright stripe \\
\texttt{vertical} & Image containing a vertical bright stripe \\
\bottomrule
\end{tabular}
\caption{Image classification dataset used in the experiment.}
\end{table}

The input of the model is an RGB image of size \(64 \times 64\), and the output is the predicted class of the input image.

\section{CNN Architecture}

The implemented network follows the VGG19 depth pattern with five convolutional blocks.
Every convolution uses a \(3 \times 3\) kernel, stride 1, and padding 1, so the spatial size is preserved inside each block.
Each convolution is followed by batch normalization and ReLU activation.
At the end of every block, max pooling with a \(2 \times 2\) kernel and stride 2 reduces the width and height by half.

The input is an RGB tensor of size \(3 \times 64 \times 64\).
The complete convolutional path is:

\begin{table}[H]
\centering
\begin{tabular}{lccc}
\toprule
Stage & Convolutions & Output channels & Spatial size after pooling \\
\midrule
Input & -- & 3 & \(64 \times 64\) \\
Block 1 & 2 & 8 & \(32 \times 32\) \\
Block 2 & 2 & 16 & \(16 \times 16\) \\
Block 3 & 4 & 32 & \(8 \times 8\) \\
Block 4 & 4 & 64 & \(4 \times 4\) \\
Block 5 & 4 & 64 & \(2 \times 2\) \\
\bottomrule
\end{tabular}
\caption{Convolutional blocks and tensor dimensions.}
\end{table}

The five blocks contain \(2+2+4+4+4=16\) convolutional layers.
For this CPU experiment, the standard VGG channel widths are multiplied by \texttt{channel\_scale = 0.125}, producing \(8,16,32,64,64\) channels.
This preserves the depth and block organization of VGG19 while reducing computation.

After the final pooling layer, adaptive average pooling keeps the feature map at \(2 \times 2\).
The resulting \(64 \times 2 \times 2=256\) features are flattened and passed through the classifier:

\begin{table}[H]
\centering
\begin{tabular}{lcc}
\toprule
Stage & Transformation & Activation \\
\midrule
Fully connected 1 & \(256 \rightarrow 128\) & ReLU and Dropout \(0.5\) \\
Fully connected 2 & \(128 \rightarrow 128\) & ReLU and Dropout \(0.5\) \\
Output & \(128 \rightarrow 2\) & Class logits \\
\bottomrule
\end{tabular}
\caption{Classifier architecture.}
\end{table}

The network therefore contains 16 convolutional layers and 3 fully connected layers, giving 19 trainable layers in the VGG19 pattern.
It has \(364714\) trainable parameters in this scaled configuration.
The final layer returns two logits, one for each class; an explicit softmax layer is unnecessary because PyTorch cross-entropy loss applies the required log-softmax operation internally.

\section{Experiment Output and Metrics}

The generated dataset is balanced and divided into independent training, validation, and test sets:

\begin{table}[H]
\centering
\begin{tabular}{lrrr}
\toprule
Split & Horizontal & Vertical & Total \\
\midrule
Training & 40 & 40 & 80 \\
Validation & 10 & 10 & 20 \\
Test & 10 & 10 & 20 \\
\bottomrule
\end{tabular}
\caption{Dataset split used in the experiment.}
\end{table}

The experiment configuration is:

\[
\text{image size} = 64 \times 64,\quad
\text{epochs} = 10,\quad
\text{batch size} = 8,\quad
\text{learning rate} = 0.001
\]

The model was initialized randomly and trained from scratch on CPU with seed 7.
Cross-entropy was used as the loss function, and Adam was used as the optimizer.

\begin{table}[H]
\centering
\begin{tabular}{ccccc}
\toprule
Epoch & Train Loss & Train Accuracy & Validation Loss & Validation Accuracy \\
\midrule
1 & 0.4859 & 0.6625 & 0.6938 & 0.5000 \\
2 & 0.0841 & 1.0000 & 1.3944 & 0.5000 \\
3 & 0.0380 & 0.9750 & 3.9379 & 0.5000 \\
4 & 0.0063 & 1.0000 & 0.4641 & 0.9500 \\
5 & 0.0041 & 1.0000 & 0.0000 & 1.0000 \\
10 & 0.0016 & 1.0000 & 0.0000 & 1.0000 \\
\bottomrule
\end{tabular}
\caption{Selected training and validation output.}
\end{table}

At epoch 1, the training accuracy was \(66.25\%\), while the validation accuracy was \(50\%\), which is the chance-level result for two balanced classes.
During epochs 2 and 3, training accuracy rapidly reached nearly \(100\%\), but validation accuracy remained at \(50\%\) and validation loss increased to \(3.9379\).
This indicates that the early model was making confident predictions on the training set before its behavior generalized to the validation set.
The small dataset and the still-developing batch-normalization statistics can make these first validation measurements unstable.

At epoch 4, validation accuracy increased to \(95\%\) and validation loss decreased to \(0.4641\).
By epoch 5, both training and validation accuracy reached \(100\%\), and validation loss fell to approximately \(2.19 \times 10^{-6}\).
Validation accuracy then remained at \(100\%\) through epoch 10.
Training loss fluctuated slightly after epoch 5 because mini-batch optimization and dropout remain stochastic during training, but training accuracy stayed at \(100\%\).

These results show that the model converged quickly after the initial unstable epochs.
However, the perfect validation result should be interpreted in the context of the small synthetic dataset: horizontal and vertical stripes have a strong visual difference, so the result does not imply the same performance on a more complex real-world image dataset.

The final test metrics are:

\begin{table}[H]
\centering
\begin{tabular}{cc}
\toprule
Metric & Value \\
\midrule
Accuracy & 1.0000 \\
Macro Precision & 1.0000 \\
Macro Recall & 1.0000 \\
Macro F1-score & 1.0000 \\
Cross-entropy loss & \(4.88 \times 10^{-6}\) \\
\bottomrule
\end{tabular}
\caption{Final classification metrics on the test set.}
\end{table}

The confusion matrix is:

\begin{table}[H]
\centering
\begin{tabular}{ccc}
\toprule
True class & Predicted horizontal & Predicted vertical \\
\midrule
horizontal & 10 & 0 \\
vertical & 0 & 10 \\
\bottomrule
\end{tabular}
\caption{Confusion matrix on the test set.}
\end{table}

The model correctly classifies all test images in this dataset.
The task is simple because the two classes are visually very different, but it is still useful for verifying that the CNN pipeline works end-to-end.

\section{Conclusion}

In this labwork, a VGG19-style convolutional neural network was implemented using PyTorch.
The model was trained from scratch without pretrained weights.
The implementation includes dataset generation, image loading, CNN training, test evaluation, and classification metrics.

The final result shows that the implemented CNN can solve the selected image classification problem successfully.
The model reaches \(1.0000\) accuracy, precision, recall, and F1-score on the test set.

\end{document}
